\section{The Research Problem}

\subsection{The Surface Problem}

A \vm{}-hosted workflow orchestration system becomes unreliable under concurrent workload
because all jobs share the same compute resources. This is an observable and known failure
mode.

\subsection{The Deeper Problem --- The Research Gap}

The surface problem is not the research gap. The gap exists at the intersection of two
bodies of literature that have not yet been connected.

Existing research evaluates Kubernetes performance at the \textit{infrastructure level} ---
measuring container overhead, pod scheduling latency, and cluster autoscaler response under
generic workloads \citep{cilic2023}. Separately, research on workflow orchestration systems
focuses on scheduling algorithms, \textsc{dag} optimisation, and system architecture. What
is missing is empirical measurement of how infrastructure-level Kubernetes mechanisms ---
pod isolation and resource limits --- actually affect \textit{application-level} workflow
execution behaviour under controlled concurrent workload conditions.

As \citet{nanda1991} demonstrate, shared execution environments degrade non-linearly under
load. As \citet{cilic2023} show, Kubernetes introduces measurable overhead. But the specific
workload threshold at which migrating from \vm{}-based to Kubernetes-based workflow
execution changes the balance between those two forces has not been empirically identified.

This thesis bridges that gap by measuring system behaviour at the pipeline execution level
rather than at the infrastructure level alone.

\subsection{The Literature Gap Statement}

Existing studies focus on Kubernetes performance at the infrastructure level, but there is
limited empirical evidence on how these mechanisms affect application-level workflow
orchestration systems under controlled concurrent workload conditions. This thesis addresses
that specific gap.
